\Needspace{3.8in}
\subsection{Segment Registers}
Segment registers define a currently addressable memory window plus a simple
address translation step. They are not the memory-protection mechanism; paging
supplies memory access permissions and final address translation. When the
segment mantissa field is zero, the segment entry is disabled. When the mantissa
is nonzero, an address passes through segmentation first and then through
paging.

The bounds-only bit enables a check-only mode for segment translation: the
segment bounds are tested, but the segment base address is not added to the
translated address.

@REGISTER_TABLE@

\begin{manuallistedformatdiagram}{Segment Register Format}{5}
\manualformatrow{segment[63:0]}{%
\manualformatfield{base\_page}{52}
\manualformatfield{e}{5}
\manualformatfield{m}{6}
\manualformatfield{b}{1}
}
\end{manuallistedformatdiagram}


@FIELD_TABLE@

\begin{align*}
\text{base byte address} &= \mathit{base\_page} \times 4096\\
\text{segment size} &= m \times 2^e \times 4096\\
\text{limit byte address} &= (\mathit{base\_page} + m \times 2^e) \times 4096\\
b=0:\quad \text{segmented address} &= \text{base byte address} + \text{offset},\quad 0 \le \text{offset} < \text{segment size}\\
b=1:\quad \text{segmented address} &= \text{offset},\quad \text{base byte address} \le \text{offset} < \text{limit byte address}
\end{align*}

The base, segment size, and limit are computed in an unsigned domain wide
enough to represent the full 64-bit linear-address space plus one. If the limit
exceeds \(2^{64}\), the segment register image is invalid and any access through
it reports \texttt{PAGE\_FAULT}.

If m is zero, segmentation is disabled for that segment selector and the
effective address is already the linear address. If m is nonzero, the segment
bounds are checked before paging. In normal mode (b = 0), the segment base is
added to the offset after the bounds check. In bounds-only mode (b = 1), the
base is not added; the segment only constrains which linear-address window may
be accessed.
